% Self-contained explanation corresponding to the Lean formalization.
\documentclass[11pt]{article}
\usepackage[T1]{fontenc}
\usepackage{amsmath,mathtools,amssymb}
\usepackage{amsthm}
\usepackage{graphicx}
\usepackage{enumitem}
\usepackage{fullpage}
\usepackage{booktabs}
\usepackage{longtable}
\usepackage[bookmarks=true,bookmarksopen=true,bookmarksnumbered=true,bookmarksdepth=2]{hyperref}
\usepackage{cleveref}

% Theorem environments.  Each kind carries its own counter and is numbered
% continuously, which is what llncs did before this became a note.
\theoremstyle{plain}
\newtheorem{theorem}{Theorem}
\newtheorem{assumption}[theorem]{Assumption}
\newtheorem{lemma}{Lemma}
\newtheorem{proposition}{Proposition}
\newtheorem{corollary}{Corollary}
\newtheorem{definition}{Definition}
\theoremstyle{remark}
\newtheorem{remark}{Remark}

\crefname{assumption}{Assumption}{Assumptions}
\Crefname{assumption}{Assumption}{Assumptions}
\crefname{appendix}{Appendix}{Appendices}
\Crefname{appendix}{Appendix}{Appendices}

\newcommand{\gaind}{\operatorname{gain}_{\delta}}
\newcommand{\gpi}{g_{\pi}}
\newcommand{\ghat}{\widehat g}
\newcommand{\gtilde}{\widetilde g}
\newcommand{\pihat}{\widehat{\pi}}
\newcommand{\sigmatilde}{\widetilde{\sigma}}
\newcommand{\Pot}{\Phi}
% Black-pebble weight on levels #1 through #2.  The dummy index is `j`, so no
% argument may itself be `j`.
\newcommand{\spend}[2]{\sum_{j=#1}^{#2}\rho_j}

\title{The Hardness of Proofs of Space Beyond Depth Robustness}
\author{ }
\date{\today}

\begin{document}
\maketitle

\Cref{sec:model} fixes the pebbling model, states the two graph hypotheses and states
the main theorem; \cref{sec:footprint-proof} proves the footprint and search lemmas it
rests on; \cref{sec:latency} assembles them into the latency bound.
\cref{app:filecoin} carries out the Filecoin-shaped specialization, and
\cref{app:lean} maps every numbered item to the Lean~4 declaration that states or proves
it.

\section{Model and main result}
\label{sec:model}
% Model and main result.
\subsection{Layered graphs and pebble configurations}

Let \(G\) have layers \(V_\ell,\ldots,V_1\), each containing \(n>0\)
nodes.  We follow Reyzin's numbering~\cite{FOCS:Reyzin24}: the top layer is
\(V_1\), the bottom layer \(V_\ell\) carries the targets, and the index \(j\) of a
layer \(V_j\) is called its \emph{level}.  Edges are acyclic and are either
intra-layer or go from \(V_{j-1}\) to \(V_j\), so the analysis reads the stack from
the challenge at level \(\ell\) upward to level \(1\).  For \(X\subseteq V_j\), let
\(\operatorname{pred}_j(X)\subseteq V_{j-1}\) be the set of inter-layer predecessors
of \(X\).  An expansion profile \(\beta:[0,1]\to[0,1]\) satisfies
\begin{equation}
 |\operatorname{pred}_j(X)|\ge
 \beta(|X|/n)n .
 \label{eq:expansion}
\end{equation}

Fix \(\pi\in(0,1)\) and \(\alpha_\pi>0\).  We assume uniform intra-layer
depth robustness in the following direct form:
\begin{equation}
 \begin{split}
  X\subseteq V_j,\quad |X|\ge\pi n
  \quad\Longrightarrow\quad
  G[X]\text{ contains a directed path with at least }\alpha_\pi n
  \text{ nodes}.
 \end{split}
 \label{eq:depth-robust}
\end{equation}
The inequalities are over the reals, so they encode the required integer rounding.

A pebble configuration consists of black sets \(\mathcal{B}_j\subseteq V_j\) and red sets
\(\mathcal{R}_j\subseteq V_j\).  Put \(\rho_j=|\mathcal{B}_j|/n\) and assume
\begin{equation}
  \spend{1}{\ell}\le\rho,
  \qquad |\mathcal{R}_j|\le\delta n\quad\text{for every }j.
  \label{eq:budgets}
\end{equation}
Black pebbles thus share one global budget; red pebbles have a per-layer budget.  A
node is \emph{unpebbled} if it belongs to neither set in its layer.

For a challenge set \(S\subseteq V_\ell\), let \(F_j(S)\) be the nodes of \(V_j\)
from which an unpebbled directed path reaches \(S\), and write
\(w_j(S)=|F_j(S)|/n\).  We require the challenge set to contain no red pebble and to satisfy
\begin{equation}
 S\cap \mathcal{R}_\ell=\varnothing,\qquad |S|/n\ge\zeta_\delta .
 \label{eq:challenge}
\end{equation}

\subsection{Expansion calculus}

Define
\[
 \beta_\delta(x):=\beta(x)-\delta,\qquad
 \gaind(x):=\beta(x)-\delta-x,\qquad
 \gpi:=\gaind(\pi),\qquad
 \bar\pi:=1-\beta(\pi).
\]
Let \(\beta\) be Reyzin's degree-\(d\) Chung expansion profile~\cite[Claim 11]{FOCS:Reyzin24}.
It maps \([0,1]\) into itself, satisfies \(\beta(0)=0\), is strictly increasing, and
strictly expands every point of \((0,1)\).  It also satisfies the reversal
identity~\cite[Claim 13]{FOCS:Reyzin24}
\begin{equation}
 \beta(1-\beta(x))=1-x\qquad(0<x<1).
 \label{eq:reversal}
\end{equation}
We assume that \(\beta\) is concave and that \(\beta(x)-x\) has a unique maximizer
\(\alpha_g<\pi\).  Finally,
\(\alpha_\delta^{\min}\le\alpha_g\le\alpha_\delta^{\max}\) are the two specified
zeros of \(\gaind\), and \(\gpi>0\).  A concave function on a segment is at least the
smaller of its endpoint values~\cite[Claim 15]{FOCS:Reyzin24}; this makes \(\gaind\)
nonnegative on \([\alpha_\delta^{\min},\alpha_\delta^{\max}]\), and it is the only way
concavity is used below.

Choose a source weight
\(\sigma\in(\alpha_\delta^{\min},\pi)\), and define the tracking quantities
\begin{equation}
 \pihat:=\min\{\bar\pi,\sigma,1-\beta(\sigma)\},
 \qquad
 \ghat:=\min\{\gpi,\gaind(\sigma)/2\}.
 \label{eq:tracking}
\end{equation}
The reversal identity gives
\(\gaind(1-\beta(\sigma))=\gaind(\sigma)\).  Hence concavity, applied between
the two endpoints \(\pihat\) and \(\pi\), yields
\begin{equation}
 \gaind(x)\ge\ghat\qquad(\pihat\le x\le\pi).
 \label{eq:tracking-gain}
\end{equation}
Both \(\ghat\) and \(\beta_\delta(\pi)-\pihat\) are positive.

The source condition also says that the gain at \(\sigma\) is \emph{twice} the tracking
constant, and concavity carries that up an initial segment.  Fix a \emph{doubled-gain
mid-point} \(\sigmatilde\) and the associated \emph{growth potential}~\(\Pot\):
\begin{equation}
 \sigmatilde\in[\sigma,\pi]
 \quad\text{with}\quad
 \gaind(\sigmatilde)\ge2\ghat ,
 \qquad
 \Pot(v):=\frac{\min\{v,\sigmatilde\}-\sigma}{2\ghat}
           +\frac{\max\{v,\sigmatilde\}-\sigmatilde}{\ghat} .
 \label{eq:midpoint}
\end{equation}
By concavity applied between \(\sigma\) and \(\sigmatilde\), \(\gaind\ge2\ghat\) on all of
\([\sigma,\sigmatilde]\).  Thus \(\Pot\) measures progress from \(\sigma\) in levels,
charging \(2\ghat\) per level below \(\sigmatilde\) and \(\ghat\) per level above it.
The choice \(\sigmatilde=\sigma\) is always admissible and gives
\(\Pot(\pi)=(\pi-\sigma)/\ghat\), which
is what the single-rate analysis charges; a genuine mid-point is what
\cref{lem:potential-window} converts into a shorter growth phase.  We use the scalar
conditions
\begin{equation}
 \alpha_\delta^{\min}<\zeta_\delta-\rho,\qquad
 \zeta_\delta\le\alpha_\delta^{\max},\qquad \rho>0.
 \label{eq:scalar-conditions}
\end{equation}
The first inequality keeps every challenge footprint in the positive-gain interval
even after the entire black budget is spent.  Together with concavity it also
makes the gain at the challenge floor \(\zeta_\delta-\rho\) positive, so
\begin{equation}
 \gtilde:=\min\{\gaind(\zeta_\delta-\rho),\gpi\}>0 .
 \label{eq:gzeta}
\end{equation}
Write \(\lceil x\rceil_+:=\max\{0,\lceil x\rceil\}\).  Three capacities drive the
theorem: the levels the searches can consume, the levels one extension attempt can
span, and the number of chain breaks.  The span has two bounds: \(h_0\) assigns the
whole budget \(\rho\) to each attempt, while \(h_1\) uses disjointness to count that
budget once across the stack.  Its growth term prices the growth phase by the better of
the single-rate count and the potential of \cref{lem:potential-window}.  In order,
\begin{align}
 s&:=\left\lceil\frac{\pi-\zeta_\delta+\rho}{\gtilde}\right\rceil_+
  +\left\lceil\frac{\rho}{\ghat}\right\rceil-1,
 \label{eq:cap-search}\\
 h_0&:=\max\left\{1,
   \left\lfloor\frac{\pi-\sigma+\rho}{\ghat}\right\rfloor\right\}
  +2\left\lceil\frac{\rho}{\ghat}\right\rceil,
 \label{eq:cap-span}\\
 h_1&:=1+\min\left\{\max\left\{1,\frac{\pi-\sigma}{\ghat}\right\},\
   \Pot(\pi)+1\right\},
 \label{eq:cap-span-global}\\
 b^{\max}&:=\left\lceil\frac{\rho}{\beta_\delta(\pi)-\pihat}\right\rceil-1.
 \label{eq:cap-break}
\end{align}
The corresponding level offsets are
\begin{equation}
 s_0:=s+b^{\max}h_0,\qquad
 s_1:=s+\frac{2\rho}{\ghat}+b^{\max}h_1 ,
 \label{eq:global-constants}
\end{equation}
and the joint ledger of \cref{lem:joint-ledger} contributes a third,
\begin{equation}
 s_2:=
 \begin{cases}
  \left(1+\dfrac{\pi-\zeta_\delta}{\gtilde}\right)_{\!+}
   +\dfrac{2\rho}{\min\{\ghat,\gtilde\}} & \text{if }b^{\max}=0,\\[2.5ex]
  \infty & \text{otherwise,}
 \end{cases}
 \label{eq:joint-constants}
\end{equation}
writing \((y)_+:=\max\{0,y\}\).  These give the chain-length bound
\begin{equation}
 z_{\min}(\ell):=\max\left\{1,
 \left\lceil\frac{\ell-s_1}{(b^{\max}+1)h_1}\right\rceil_+,
 \left\lceil\frac{\ell-s_2}{h_1}\right\rceil_+,
 \left\lfloor\frac{\ell-s_0}{(b^{\max}+1)h_0}\right\rfloor+1\right\}.
 \label{eq:zmin}
\end{equation}
The third entry vanishes exactly when a break can be paid for, which is what
\(s_2=\infty\) records: the joint ledger needs the searches and the attempts to occupy
disjoint ranges of levels, and a restart search destroys that.
All denominators are positive, so \(s\ge0\), \(h_0\ge3\),
\(b^{\max}\ge0\), \(h_1\ge2\), \(s_2>0\) and \(s_0\ge0\).
The four entries of \eqref{eq:zmin} are independent lower bounds on the chain length,
so taking their maximum loses nothing: \(s_1\) and \(s_2\) are not comparable
in general, and neither dominates for every parameter choice.  Every quantity in
\eqref{eq:gzeta}--\eqref{eq:zmin} is a function of
\(\beta,\delta,\pi,\rho,\sigma,\sigmatilde,\zeta_\delta\) alone; only \(z_{\min}\) also reads
\(\ell\).  The three capacities are proved in
\cref{lem:infertile-capacity,lem:first-source} for \(s\), in
\cref{lem:extension-attempt} for \(h_0\), and in \cref{lem:break-charge} for
\(b^{\max}\); the joint constants in \cref{lem:joint-ledger};
the global span \(h_1\) is established inside
\cref{lem:global-ledger}.

\begin{theorem}[Latency theorem]
\label{thm:latency}
Assume \eqref{eq:expansion}--\eqref{eq:scalar-conditions} and
\(\sigma<\alpha_\pi\).  If \(\ell>s_0\), every pebble configuration satisfying
\eqref{eq:budgets} and every challenge set \(S\) satisfying \eqref{eq:challenge} have
an unpebbled directed path \(Q\), contained in the footprint of \(S\) and ending in
\(S\), such that
\begin{equation}
 |Q|\ge \alpha_\pi n+
 \bigl(z_{\min}(\ell)-1\bigr)
 (\alpha_\pi-\sigma)n.
 \label{eq:main-path-length}
\end{equation}
In particular, for fixed parameters the path length is \(\Omega(\ell n)\).
\end{theorem}

The proof first works with footprint bounds and then combines the resulting graph
paths.

\section{Footprint growth and search}
\label{sec:footprint-proof}
% Footprint growth and search.
\subsection{Footprint bounds}

The \emph{footprint bound} started at level \(t\) with weight \(y\) is given by
Reyzin's recurrence~\cite{FOCS:Reyzin24}:
\begin{equation}
 f_t=y,\qquad
 f_{j-1}=\max\{0,\beta_\delta(f_j)-\rho_{j-1}\}\qquad(j\le t).
 \label{eq:footprint-bound}
\end{equation}
For the challenge we instead initialize at the bottom level,
\begin{equation}
 f_\ell=\max\{0,\zeta_\delta-\rho_\ell\}.
 \label{eq:challenge-bound}
\end{equation}
Both run from their starting level upward, toward level \(1\).

\subsection{The challenge floor and the first chain source}

Call a challenge level \(j\) \emph{fertile} when \(f_j\ge\pi\).  For \(g>0\), call
a level \(t\) \emph{\(g\)-expandable} when
\begin{equation}
 \spend{t-i}{t-1}\le(i+1)g
 \qquad\text{for every }i\ge1.
 \label{eq:expandable}
\end{equation}
This constrains the black weight in every window of levels lying above \(t\).

\begin{lemma}[General challenge floor]
\label{lem:challenge-floor}
Under \eqref{eq:scalar-conditions}, the challenge footprint bound satisfies, at
every level \(m\),
\begin{equation}
 \zeta_\delta-\spend{m}{\ell}\le f_m\le
 \alpha_\delta^{\max}.
 \label{eq:challenge-invariants}
\end{equation}
Consequently \(f_m\ge\zeta_\delta-\rho\), every gain
\(\gaind(f_m)\) is nonnegative, and every infertile level has gain at least the
constant \(\gtilde>0\) of \eqref{eq:gzeta}.
\end{lemma}

\begin{proof}
We prove both inequalities in \eqref{eq:challenge-invariants}
simultaneously~\cite[Claims 4 and 5]{FOCS:Reyzin24}; the same induction also gives the
\(\gtilde\) floor.  They hold at \(m=\ell\) by \eqref{eq:challenge-bound}, the
nonnegativity of \(\zeta_\delta\), and
\(\zeta_\delta\le\alpha_\delta^{\max}\).  Suppose they hold at \(m\).  Since
\(\spend{m}{\ell}\le\rho\) by \eqref{eq:budgets}, the lower bound gives
\[
 f_m\ge\zeta_\delta-\rho>\alpha_\delta^{\min}.
\]
Together with the upper bound, this places \(f_m\) in the active interval, so
\(\gaind(f_m)\ge0\).  The recurrence now advances the lower invariant.

For the upper invariant, \(\beta_\delta\) is strictly increasing and fixes
\(\alpha_\delta^{\max}\), because
\(\gaind(\alpha_\delta^{\max})=0\).  Hence
\(\beta_\delta(f_m)\le\alpha_\delta^{\max}\); subtracting nonnegative black weight
and truncating at zero preserves this upper bound.  This completes the induction.

The floor \(f_m\ge\zeta_\delta-\rho\) follows again from the total budget.  If
\(f_m<\pi\), then \(f_m\in[\zeta_\delta-\rho,\pi]\).  Concavity of
\(\gaind\) implies that its value on this interval is at least the minimum of
its endpoint values, namely \(\gtilde\).
\end{proof}

\begin{lemma}[Infertile capacity]
\label{lem:infertile-capacity}
The challenge footprint bound has at most
\[
 \left\lceil\frac{\pi-\zeta_\delta+\rho}{\gtilde}\right\rceil_+
\]
infertile levels in total.  This is the first summand of \(s\) in
\eqref{eq:cap-search}.
\end{lemma}

\begin{proof}
Suppose there are \(k>0\) infertile levels, and let \(m\) be the smallest of them.
The other \(k-1\) lie among the levels \(m+1,\ldots,\ell\), each contributes gain
at least \(\gtilde\), and all remaining gains there are nonnegative by
\cref{lem:challenge-floor}.  The accumulated bound~\cite[Claim 3]{FOCS:Reyzin24},
applied at \(m\) with \(f_\ell\ge\zeta_\delta-\rho_\ell\) and \(f_m<\pi\), gives
\[
 \zeta_\delta-\spend{m}{\ell}+(k-1)\gtilde<\pi.
\]
Since \(\spend{m}{\ell}\le\rho\),
\[
 (k-1)\gtilde<\pi-\zeta_\delta+\rho.
\]
If \(\zeta_\delta-\rho\ge\pi\), the floor from
\cref{lem:challenge-floor} rules out an infertile level, and the displayed bound is
\(0\).  Otherwise, dividing by \(\gtilde>0\) and applying the ceiling gives the
bound.  The case \(k=0\) is immediate.
\end{proof}

\Cref{lem:infertile-capacity} replaces the black weight actually spent between the
deepest infertile level and the challenge by the whole budget \(\rho\).  The joint
ledger of \cref{lem:joint-ledger} needs the same count charged to the levels it is
levied on, so we record that form as well.

\begin{lemma}[Infertile capacity in a window]
\label{lem:infertile-window}
Let \(1\le p<\ell\).  If \(k\ge1\) of the levels \(p+1,\ldots,\ell\) are
infertile for the challenge footprint bound, then
\begin{equation}
 (k-1)\gtilde<\pi-\zeta_\delta+\spend{p+1}{\ell} .
 \label{eq:infertile-window}
\end{equation}
Consequently, for every \(k\ge0\),
\begin{equation}
 k\le\left(1+\frac{\pi-\zeta_\delta}{\gtilde}\right)_{\!+}
   +\frac{1}{\gtilde}\spend{p+1}{\ell},
 \label{eq:infertile-window-head}
\end{equation}
whose first summand is the head of \(s_2\) in \eqref{eq:joint-constants}.
\end{lemma}

\begin{proof}
This is the proof of \cref{lem:infertile-capacity} with the last step omitted.  Let
\(m\ge p+1\) be the smallest infertile level among \(p+1,\ldots,\ell\); the other
\(k-1\) lie among \(m+1,\ldots,\ell\), each contributes gain at least \(\gtilde\),
and all remaining gains there are nonnegative by \cref{lem:challenge-floor}.  The
accumulated bound~\cite[Claim 3]{FOCS:Reyzin24} applied at \(m\), with
\(f_\ell\ge\zeta_\delta-\rho_\ell\) and \(f_m<\pi\), gives
\[
 \zeta_\delta-\spend{m}{\ell}+(k-1)\gtilde<\pi ,
\]
and \(\spend{m}{\ell}\le\spend{p+1}{\ell}\) because \(m\ge p+1\) and the black
weights are nonnegative.  This is \eqref{eq:infertile-window}.  Dividing by
\(\gtilde>0\) gives
\(k\le1+(\pi-\zeta_\delta)/\gtilde+\spend{p+1}{\ell}/\gtilde\), and the first two
terms are at most their positive part.  When \(k=0\),
\eqref{eq:infertile-window-head} holds because both terms on its right are
nonnegative.
\end{proof}

The same search runs in three places, on two different footprint bounds: at the start
of the chain, after every break, and inside every extension attempt.  Starting at a
level \(t\) and going up, it looks for a level that is both fertile and
\(\ghat\)-expandable.  At a level \(p\) it stops if \(p\) is such a level; if \(p\) is
infertile it moves to \(p-1\); and if \(p\) is fertile but not \(\ghat\)-expandable,
there is an \(i\ge1\) with
\[
 \spend{p-i}{p-1}>(i+1)\ghat,
\]
and it skips that \emph{blocked} window, moving to \(p-i-1\).  Blocked windows
skipped this way occupy disjoint ranges of levels, so their black weight is charged
once.  The search may also stop because it has gone past the top of the graph, or past
a prescribed cutoff.  Nothing in the search uses what fertility means, so we may run it
on any footprint bound.  Running it on the challenge bound, and counting infertile
levels with \cref{lem:infertile-capacity}, gives the first structural event.

\begin{lemma}[First fertile expandable level]
\label{lem:first-source}
If \(\ell>s\), the search on the challenge footprint bound started at level \(\ell\)
stops at a level \(b\ge\ell-s\) that is both fertile and \(\ghat\)-expandable.
\end{lemma}

\begin{proof}
Let \(u\) count the levels the search skips as infertile and \(v\) those it skips
inside blocked windows; at every point of its run it stands at level \(\ell-u-v\).
The infertile skips are at distinct infertile levels, so \(u\) is at most the count of
\cref{lem:infertile-capacity}.  If \(v>0\), summing the strict inequalities of
the skipped windows over their disjoint ranges gives \(v\ghat<\rho\), whence
\(v\le\lceil\rho/\ghat\rceil-1\).  The two bounds add to \(s\) by
\eqref{eq:cap-search}, so the search never goes past level \(\ell-s\ge1\).  It
therefore stops inside the graph, at a fertile \(\ghat\)-expandable level
\(b\ge\ell-s\).
\end{proof}

The challenge footprint at \(b\) has weight at least \(\pi\).  Condition
\eqref{eq:depth-robust} therefore gives an unpebbled
path of length at least \(\alpha_\pi n\) inside that footprint.  Since
\(\sigma<\alpha_\pi\), take the first \(\lceil\sigma n\rceil\) nodes of that
path as a source set.  Each source node begins a suffix that reaches \(S\), and the
source has normalized weight at least \(\sigma\).  This source set is the first
\emph{chain source}.

\subsection{Tracking one source}

Consider the footprint bound of a chain source, initialized by \(f_t=\sigma\) at an
expandable level \(t\).  The next lemma is the analytic reason for
\eqref{eq:tracking}.

\begin{lemma}[Mirror floor]
\label{lem:mirror-floor}
Let \(t\) be \(\ghat\)-expandable and let \(m\le t\).  If the source footprint
bound is at most \(\pi\) at every level of \(m,\ldots,t\), then it is at least
\(\pihat\) at every level of \(m,\ldots,t\) --- indeed at least \(\sigma\).
\end{lemma}

\begin{proof}
The claim at \(m=t\) follows from \(\pihat\le\sigma\).  Proceed by strong induction
downward.  At the source, \(\gaind(\sigma)\ge2\ghat\).  At each of the \(t-m-1\)
levels strictly between \(t\) and \(m\), the induction hypothesis and the upper
assumption place the bound in \([\pihat,\pi]\), where \eqref{eq:tracking-gain} gives
gain at least \(\ghat\).  Thus the gain sum in the accumulated
bound~\cite[Claim 3]{FOCS:Reyzin24} at \(m\) is at least
\((t-m+1)\ghat\).  Expandability with \(i=t-m\) bounds
\(\spend{m}{t-1}\) by the same quantity, and consequently
\(f_m\ge f_t=\sigma\ge\pihat\).
\end{proof}

\begin{lemma}[Growth window]
\label{lem:growth-window}
Suppose \(t_1<t\) is the first level above \(t\) at which the source footprint bound
is larger than \(\pi\), and put
\(\rho_{\mathrm{grow}}=\spend{t_1+1}{t-1}\).  Then
\begin{equation}
 t-t_1\le
 \max\left\{1,
 \left\lfloor\frac{\pi-\sigma+\rho_{\mathrm{grow}}}{\ghat}\right\rfloor\right\}.
 \label{eq:growth-window}
\end{equation}
\end{lemma}

\begin{proof}
If \(t_1=t-1\), the claim follows from the explicit \(1\) in the maximum.
Otherwise apply \cref{lem:mirror-floor} through level \(t_1+1\).  The source gain is
at least \(2\ghat\), and every later pre-crossing gain is at least \(\ghat\).  Since
\(f_{t_1+1}\le\pi\), the accumulated bound~\cite[Claim 3]{FOCS:Reyzin24} gives
\[
 (t-t_1)\ghat\le\pi-\sigma+\rho_{\mathrm{grow}}.
\]
The left side is an integer multiple of \(\ghat\), so taking the floor proves the
claim.
\end{proof}

\Cref{lem:growth-window} charges every level of the growth phase at the single rate
\(\ghat\), which lower-bounds the gain throughout the tracking interval
\([\pihat,\pi]\).  In between the gain can be larger, and \eqref{eq:midpoint} certifies
how much larger on an initial segment.  A potential converts that into a level count.

\begin{lemma}[Two-piece growth window]
\label{lem:potential-window}
In the situation of \cref{lem:growth-window},
\begin{equation}
 t-t_1\le\Pot(\pi)+1+\frac{\rho_{\mathrm{grow}}}{\ghat},
 \label{eq:potential-window}
\end{equation}
with \(\Pot\) the growth potential of \eqref{eq:midpoint}.
\end{lemma}

\begin{proof}
The potential is nondecreasing, \(\Pot(\sigma)=0\), and it is
\(1/\ghat\)-Lipschitz: its slope is \(1/(2\ghat)\) below \(\sigmatilde\) and \(1/\ghat\)
above, and \(1/\ghat\) is the larger of the two.

\emph{One free level advances \(\Pot\) by at least \(1\).}  Let
\(\sigma\le v\le\pi\).  If \(v\le\sigmatilde\) then \(\gaind(v)\ge2\ghat\) by
\eqref{eq:midpoint} and concavity, so
\(\beta_\delta(v)-v\ge2\ghat\); the increment
\(\Pot(\beta_\delta(v))-\Pot(v)\) is at least
\((\beta_\delta(v)-v)/(2\ghat)\ge1\), because the slope of \(\Pot\) on
\([v,\beta_\delta(v)]\) is at least \(1/(2\ghat)\) throughout, whether or not that
interval crosses \(\sigmatilde\).  If \(v\ge\sigmatilde\) then \eqref{eq:tracking-gain} gives
\(\beta_\delta(v)-v\ge\ghat\), and the slope of \(\Pot\) is exactly
\(1/\ghat\) there, so the increment is again at least \(1\).

\emph{Telescoping.}  Put \(m_i=t-i\) and \(N=t-t_1-1\), so that
\(f_{m_i}\le\pi\) for \(0\le i\le N\).  By \cref{lem:mirror-floor} each such
\(f_{m_i}\) is at least \(\sigma\), so the previous paragraph applies at every step.
The recurrence \eqref{eq:footprint-bound} gives
\(f_{m_{i+1}}\ge\beta_\delta(f_{m_i})-\rho_{m_{i+1}}\), and monotonicity together
with the Lipschitz bound yields
\[
 \Pot(f_{m_{i+1}})\ge\Pot(\beta_\delta(f_{m_i}))-\frac{\rho_{m_{i+1}}}{\ghat}
 \ge\Pot(f_{m_i})+1-\frac{\rho_{m_{i+1}}}{\ghat} .
\]
Summing from \(i=0\) to \(N-1\), and using \(\Pot(f_t)=\Pot(\sigma)=0\),
\[
 \Pot(f_{t_1+1})\ge N-\frac{1}{\ghat}\spend{t_1+1}{t-1}
 =N-\frac{\rho_{\mathrm{grow}}}{\ghat} .
\]
Since \(f_{t_1+1}\le\pi\) and \(\Pot\) is nondecreasing, the left side is at most
\(\Pot(\pi)\), so
\(t-t_1=N+1\le\Pot(\pi)+1+\rho_{\mathrm{grow}}/\ghat\).
\end{proof}

Comparing the two window bounds is a one-line computation:
\begin{equation}
 \Pot(\pi)+\frac{\sigmatilde-\sigma}{2\ghat}=\frac{\pi-\sigma}{\ghat},
 \label{eq:potential-gap}
\end{equation}
so the potential saves \((\sigmatilde-\sigma)/(2\ghat)\) levels and pays back the one free level
that \cref{lem:mirror-floor} extracts at the source.  It therefore wins exactly when
\(\sigmatilde-\sigma>2\ghat\), and the growth term \(h_1-1\) of \eqref{eq:cap-span-global}
takes whichever bound is smaller.  Both \cref{lem:growth-window,lem:potential-window}
hold, so
\begin{equation}
 t-t_1\le h_1-1+\frac{\rho_{\mathrm{grow}}}{\ghat}
 \label{eq:growth-constant-window}
\end{equation}
in either case; this is the only form the accounting uses.

Once a footprint bound is fertile, the same search has a local bound depending only
on the black weight inside its window.

\begin{lemma}[Fertile continuation]
\label{lem:fertile-continuation}
Let \(f_{t_1}\ge\pi\), and suppose that down to a cutoff \(e\) the bound stays in
\([\pihat,\alpha_\delta^{\max}]\).  The search beginning at \(t_1\) ends at a
level \(t_2\), either by leaving \((e,t_1]\) or by finding a fertile
\(\ghat\)-expandable level, and
\begin{equation}
 t_1-t_2\le
 \max\left\{0,
 2\left\lceil\frac{\rho_{\mathrm{cont}}}{\ghat}\right\rceil-1\right\},
 \qquad \rho_{\mathrm{cont}}:=\spend{t_2+1}{t_1-1}.
 \label{eq:continuation-window}
\end{equation}
\end{lemma}

\begin{proof}
Run the same search, now calling a level fertile when the tracked bound is at least
\(\pi\) there, and stopping also when it reaches \(e\).  Let \(u\) and \(v\) denote its
infertile and blocked charges.  All gains in the search are nonnegative.  At an
infertile level the value lies in \([\pihat,\pi]\), so its gain is at least
\(\ghat\).  If \(u>0\), apply the accumulated
bound~\cite[Claim 3]{FOCS:Reyzin24} at the last infertile level the search meets.  The
initial term \(f_{t_1}-\pi\) is nonnegative, so
\((u-1)\ghat<\rho_{\mathrm{cont}}\), giving
\(u\le\lceil \rho_{\mathrm{cont}}/\ghat\rceil\).  The disjoint blocked ranges satisfy
\(v\ghat<\rho_{\mathrm{cont}}\) when \(v>0\).  If \(\rho_{\mathrm{cont}}=0\),
these strict inequalities force
\(u=v=0\).  Otherwise
\(v\le\lceil \rho_{\mathrm{cont}}/\ghat\rceil-1\).  Since the search goes up exactly through the
levels counted by \(u+v\), in the positive case
\[
 t_1-t_2\le u+v\le
 2\left\lceil\frac{\rho_{\mathrm{cont}}}{\ghat}\right\rceil-1.
\]
This is the stated bound in both cases.
\end{proof}

The cutoff in \cref{lem:fertile-continuation} is the first level at which the tracked
bound drops below \(\pihat\), or the top of the graph if that never happens.
The former event has a fixed price.

\begin{lemma}[Break charge]
\label{lem:break-charge}
Suppose \(f_{t_1}\ge\pi\), \(m<t_1\), all intermediate gains are nonnegative, and
\(f_m<\pihat\).  Then
\begin{equation}
 \beta_\delta(\pi)-\pihat
 <\spend{m}{t_1-1}.
 \label{eq:break-charge}
\end{equation}
\end{lemma}

\begin{proof}
Strict monotonicity of \(\beta_\delta\) and the first recurrence step give
\[
 f_{t_1-1}\ge\beta_\delta(f_{t_1})-\rho_{t_1-1}
 \ge\beta_\delta(\pi)-\rho_{t_1-1}.
\]
Apply the accumulated bound~\cite[Claim 3]{FOCS:Reyzin24} from \(t_1-1\) down to
\(m\).  Discarding the nonnegative intermediate gains yields
\[
 f_m\ge\beta_\delta(\pi)-\spend{m}{t_1-1}.
\]
Combining this with \(f_m<\pihat\) proves \eqref{eq:break-charge}.
\end{proof}

The two phases together form one extension attempt.

\begin{lemma}[One extension attempt]
\label{lem:extension-attempt}
An extension attempt from any \(\ghat\)-expandable chain source at level \(t\) ends
at the top of the graph, at the next chain source, or at a break, and its terminal level
\(m\) satisfies
\begin{equation}
 t-m\le h_0-1,
 \label{eq:attempt-bound}
\end{equation}
with \(h_0\) the integer \eqref{eq:cap-span}.
In the break outcome, its window also satisfies \eqref{eq:break-charge}.
\end{lemma}

\begin{proof}
Follow the source footprint bound until it first exceeds \(\pi\), or until the graph
is exhausted.  The proof of \cref{lem:growth-window}, with the terminal level in
place of \(t_1\) in the latter case, bounds the growth phase by
\(\max\{1,\lfloor(\pi-\sigma+\rho)/\ghat\rfloor\}\), and
\eqref{eq:growth-constant-window} bounds it by
\(h_1-1+\rho_{\mathrm{grow}}/\ghat\).
If a fertile level is reached, run the search of
\cref{lem:fertile-continuation} until it finds a fertile expandable level, reaches
the top, or crosses the tracking floor.  These are exactly the three stated outcomes.
The continuation phase has length at most \(2\lceil\rho/\ghat\rceil-1\)
by \cref{lem:fertile-continuation} and the global budget.  Adding the two phase
bounds gives \eqref{eq:attempt-bound}.  (If the graph ends during the growth phase,
the continuation contribution is zero.)
If the tracking floor is crossed, \cref{lem:break-charge} supplies the additional
charge.  Finally, at a fertile expandable outcome, the corresponding source footprint
has weight at least \(\pi\), so \eqref{eq:depth-robust} constructs the next chain source
exactly as after \cref{lem:first-source}.
\end{proof}

\section{The latency theorem}
\label{sec:latency}
% Latency theorem.
\subsection{Building and counting chains}

A \emph{chain source} is a source set of weight at least \(\sigma\) at an expandable
level such that every node in the set has an unpebbled path to \(S\).
A \emph{chain segment} begins at a fertile expandable level found by searching the
challenge footprint bound; an unpebbled path in the corresponding footprint
supplies the first chain source.  Starting from the last chain source, repeatedly
apply \cref{lem:extension-attempt}.  A successful attempt finds the next chain source
of the same segment.  At a break, abandon that segment and resume the search on
the challenge footprint bound from the break level.  All searches and attempts go up
the stack, so their charged intervals are pairwise disjoint.

Let \(z\) be the largest number of chain sources in any completed segment, \(a\) the number
of extension attempts, and \(b\) the number of breaks.

\begin{lemma}[Counting sources and levels]
\label{lem:global-ledger}
Run the construction from level \(\ell\) and suppose that its initial search finds the
first chain source.  It then terminates with
\(z\ge1\) and natural numbers \(a,b\) satisfying
\begin{align}
 1+a&\le(b+1)(z+1),
 \label{eq:segment-ledger}\\
 \ell&\le s+a(h_0-1),
 \label{eq:level-ledger}\\
 \ell&\le s+\frac{2\rho}{\ghat}+a h_1,
 \label{eq:global-level-ledger}\\
 b&\le b^{\max} .
 \label{eq:break-ledger}
\end{align}
One of the resulting segments contains \(z\) chain sources, each obtained from
unpebbled paths in the corresponding footprints.  If \(b=0\), that segment is the only one and its
last chain source lies at a level \(m\)
with
\begin{equation}
 m\le h_0-1 .
 \label{eq:last-source-depth}
\end{equation}
\end{lemma}

\begin{proof}
Every level of the graph is charged either to a search or to an extension attempt.
Write \(u\) for the number of one-step advances past an infertile level and \(v\)
for the number of levels charged to blocked search windows, so the searches consume
\(u+v\) levels in total.  An extension attempt increases \(a\) by one and, by
\cref{lem:extension-attempt}, advances at most \(h_0-1\) levels.  Hence
\(\ell\le u+v+a(h_0-1)\).

Now \(u+v\le s\).  Restart searches operate on disjoint level ranges, so their
infertile moves are distinct infertile levels of the single challenge footprint bound,
and
\cref{lem:infertile-capacity} bounds \(u\) by the first summand of
\eqref{eq:cap-search}.  The blocked windows are disjoint as well; if their total length is
\(v>0\), summing their defining strict inequalities gives
\[
 v\ghat<\sum_{j\text{ in blocked windows}}\rho_j\le\rho
 \le\left\lceil\frac{\rho}{\ghat}\right\rceil\ghat ,
\]
so \(v\le\lceil\rho/\ghat\rceil-1\), which holds trivially when \(v=0\).  Adding
the two bounds gives \eqref{eq:level-ledger}.

For the global version, do not replace the spend of each attempt by the whole budget.
The growth phase of an attempt that sees black weight \(\rho_{\mathrm{grow}}\) spans at most
\(h_1-1+\rho_{\mathrm{grow}}/\ghat\) levels, and its continuation phase, with weight
\(\rho_{\mathrm{cont}}\), spans at most \(1+2\rho_{\mathrm{cont}}/\ghat\).  The
attempt windows are disjoint, so summing gives
\[
 \ell\le u+v+a h_1
       +\frac{1}{\ghat}\sum \rho_{\mathrm{grow}}
       +\frac{2}{\ghat}\sum \rho_{\mathrm{cont}}
 \le s+a h_1+\frac{2\rho}{\ghat},
\]
which is \eqref{eq:global-level-ledger}.

Initially there is one chain source.  A successful attempt adds one chain source to the
current segment.  A terminal or broken attempt contributes no chain source, and a break
increases the number of segments by one.  If the numbers of chain sources in the
segments are \(z_1,\ldots,z_{b+1}\), then
\[
 1+a\le\sum_{j=1}^{b+1}(z_j+1)\le(b+1)(z+1),
\]
which proves \eqref{eq:segment-ledger}.  The same induction retains the footprints and
unpebbled paths used to obtain each new chain source, including every chain source of a
segment attaining \(z\).

For \eqref{eq:break-ledger}, every break window spends strictly more than
\(\beta_\delta(\pi)-\pihat\) by \cref{lem:break-charge}, and break windows are
disjoint.  If \(b>0\), summing over them gives
\(b\bigl(\beta_\delta(\pi)-\pihat\bigr)<\rho\le
(b^{\max}+1)\bigl(\beta_\delta(\pi)-\pihat\bigr)\), so \(b\le b^{\max}\); this
holds trivially when \(b=0\).

Finally, every nonterminal attempt either finds the next chain source or causes a restart, and a
search that leaves the graph is terminal.  Since the level strictly decreases, the
construction terminates after finitely many events.

For \eqref{eq:last-source-depth}, suppose \(b=0\).  There is then one segment, and the
construction can only stop by leaving the graph, which an attempt does.  That attempt
starts at the last chain source at level \(m\) and ends at a terminal level
\(m'\le0\); by \eqref{eq:attempt-bound}, \(m-m'\le h_0-1\), and \(m'\le0\) gives
\(m\le h_0-1\).
\end{proof}

\Cref{lem:global-ledger} charges the black budget \(\rho\) three separate times.  It
appears in the infertile-capacity term of \(s\) through \cref{lem:infertile-capacity},
again in the blocked-window term \(\lceil\rho/\ghat\rceil-1\) of \(s\), and twice more
in the \(2\rho/\ghat\) of \eqref{eq:global-level-ledger}.  Yet the levels those three
charges are levied on are pairwise disjoint --- a level belongs either to a search or to
an attempt, and the construction moves strictly upward --- so one interval sum bounds all
three, and \eqref{eq:budgets} caps that sum by \(\rho\).  Collecting them is the joint
ledger.

The obstacle in general is a break: the restart search that follows it runs above levels
already charged to attempts, and the disjointness fails.  When \(b^{\max}=0\) no break
can be paid for, there is exactly one search, and it precedes every attempt.

\begin{lemma}[Joint ledger]
\label{lem:joint-ledger}
Suppose \(b^{\max}=0\) and \(\ell>s\).  Then the construction terminates with a single
segment, and its numbers \(z\ge1\) and \(a\) satisfy \eqref{eq:segment-ledger} and
\begin{equation}
 \ell\le s_2+a\,h_1 .
 \label{eq:joint-level-ledger}
\end{equation}
\end{lemma}

\begin{proof}
By \cref{lem:break-charge} a break costs more than \(\beta_\delta(\pi)-\pihat\) of
black weight, and \(b^{\max}=0\) says \(\rho\) is at most that, so no attempt breaks.
Hence there is one segment, and the only search is the initial one; it starts at level
\(\ell\) and, by \cref{lem:first-source}, stops at the first chain source
\(b_1\ge\ell-s\).  Every attempt window lies in \([1,b_1-1]\): the first attempt starts at
\(b_1\) and charges only the levels strictly above it, and every later attempt starts
higher still, while the search occupied \(b_1+1,\ldots,\ell\).  Write
\[
 \rho_{\mathrm{search}}:=\spend{b_1+1}{\ell},\qquad
 \rho_{\mathrm{attempts}}:=\spend{1}{b_1-1},
\]
so that \(\rho_{\mathrm{search}}+\rho_{\mathrm{attempts}}\le\rho\) by
\eqref{eq:budgets}.

\emph{The search.}  Let \(u\) and \(v\) be as in the proof of
\cref{lem:first-source}, so the search consumes \(u+v\) levels and
\(\ell-b_1=u+v\).  Its infertile skips are at distinct infertile levels of
\(b_1+1,\ldots,\ell\), so \eqref{eq:infertile-window-head} with \(p=b_1\) bounds \(u\) by the head of \(s_2\)
plus \(\rho_{\mathrm{search}}/\gtilde\).  Its blocked windows are disjoint subranges
of the same levels, so \(v\ghat<\rho_{\mathrm{search}}\) when \(v>0\), and
\(v\le \rho_{\mathrm{search}}/\ghat\) in either case.  Writing \(\varepsilon\) for
the head of \(s_2\) in
\eqref{eq:joint-constants} and bounding both rates below by
\(\min\{\ghat,\gtilde\}\),
\begin{equation}
 \ell-b_1\le\varepsilon
   +\frac{\rho_{\mathrm{search}}}{\gtilde}
   +\frac{\rho_{\mathrm{search}}}{\ghat}
 \le\varepsilon
   +\frac{2\rho_{\mathrm{search}}}{\min\{\ghat,\gtilde\}} .
 \label{eq:joint-search}
\end{equation}

\emph{The attempts.}  As in the proof of \eqref{eq:global-level-ledger}, an attempt
whose growth phase sees black weight \(\rho_{\mathrm{grow}}\) and whose continuation
sees \(\rho_{\mathrm{cont}}\) spans at most
\(h_1-1+\rho_{\mathrm{grow}}/\ghat\) levels growing, by
\eqref{eq:growth-constant-window}, and at most
\(1+2\rho_{\mathrm{cont}}/\ghat\) continuing, by
\cref{lem:fertile-continuation}; that is
\(h_1+\rho_{\mathrm{grow}}/\ghat+2\rho_{\mathrm{cont}}/\ghat\) in total.  The
attempt windows are disjoint, so summing over the \(a\) attempts,
\begin{equation}
 b_1\le a\,h_1
   +\frac{2\rho_{\mathrm{attempts}}}{\min\{\ghat,\gtilde\}} .
 \label{eq:joint-attempts}
\end{equation}

Adding \eqref{eq:joint-search} and \eqref{eq:joint-attempts} and using
\(\rho_{\mathrm{search}}+\rho_{\mathrm{attempts}}\le\rho\) gives
\(\ell\le\varepsilon+a\,h_1+2\rho/\min\{\ghat,\gtilde\}\), which is
\eqref{eq:joint-level-ledger} by \eqref{eq:joint-constants}.  The segment
structure \eqref{eq:segment-ledger} is unchanged, and \(z\ge1\) because
\cref{lem:first-source} produced a chain source.
\end{proof}

The four inequalities now give the chain length.

\begin{lemma}[Chain length]
\label{lem:chain-length}
Whenever \(\ell>s_0\), the construction finds a segment with
\begin{equation}
 z\ge z_{\min}(\ell)
 \label{eq:chain-length}
\end{equation}
chain sources.
\end{lemma}

\begin{proof}
If \(\ell>s_0\), then in particular \(\ell>s\), because
\(s_0=s+b^{\max} h_0\ge s\).  \Cref{lem:first-source} therefore finds the
first chain source before the graph ends, and we may apply
\cref{lem:global-ledger}.  From \eqref{eq:segment-ledger} and
\eqref{eq:break-ledger},
\[
 a\le(b^{\max}+1)(z+1)-1=(b^{\max}+1)z+b^{\max} .
\]
Substituting this into \eqref{eq:level-ledger}, and using \(z\ge1\) to make the
last step strict, gives
\begin{align*}
 \ell
 &\le s+\bigl((b^{\max}+1)z+b^{\max}\bigr)(h_0-1)\\
 &<s+\bigl((b^{\max}+1)z+b^{\max}\bigr)h_0\\
 &=s+b^{\max} h_0+z(b^{\max}+1)h_0
   =s_0+z(b^{\max}+1)h_0.
\end{align*}
Thus the constant-charge entry of \eqref{eq:zmin} is at most \(z\).

Substitution in \eqref{eq:global-level-ledger} instead gives
\[
 \ell\le s+\frac{2\rho}{\ghat}+
 \bigl((b^{\max}+1)z+b^{\max}\bigr)h_1
 =s_1+z(b^{\max}+1)h_1,
\]
so the nonnegative-ceiling entry of \eqref{eq:zmin} is also at most \(z\).

For the third entry of \eqref{eq:zmin} there is nothing to prove unless
\(b^{\max}=0\), since \(s_2=\infty\) otherwise.  When \(b^{\max}=0\),
\eqref{eq:segment-ledger} reads \(1+a\le z+1\) and \cref{lem:joint-ledger} gives
\(\ell\le s_2+a\,h_1\le s_2+z\,h_1\), so that entry is at most \(z\) as well.
Finally \(z\ge1\), proving \eqref{eq:chain-length}.
\end{proof}

\subsection{Combining the paths}

\begin{lemma}[Path length from a chain segment]
\label{lem:path-payoff}
A segment of \(z\ge1\) chain sources yields an unpebbled path in the footprint of \(S\)
of length at least
\begin{equation}
 p(z):=\alpha_\pi n+(z-1)(\alpha_\pi-\sigma)n.
 \label{eq:segment-payoff}
\end{equation}
\end{lemma}

\begin{proof}
For the first chain source, depth robustness gives a path \(Q\) of length at least
\(\alpha_\pi n\) in the challenge footprint.  The last node of \(Q\) has an
unpebbled path to \(S\); appending it cannot shorten \(Q\).  This proves
\eqref{eq:segment-payoff} for \(z=1\).

We maintain the stronger invariant that every node in the
\(z\)th chain source has an unpebbled path to \(S\) of length at least
\(z(\alpha_\pi-\sigma)n\).  It holds for the first chain source: if the source is the first
\(k=\lceil\sigma n\rceil\) nodes of \(Q\), then a source position \(i<k\)
satisfies \(i<\sigma n\), and its suffix contains at least
\[
 |Q|-i\ge\alpha_\pi n-\sigma n
       =(\alpha_\pi-\sigma)n
\]
nodes.

Suppose the invariant holds for the \(z\)th chain source \(T\), and a later
fertile footprint \(F_j(T)\) is used to obtain the next chain source.  Let \(Q\) be
the path
given by depth robustness.
Its last node reaches some \(y\in T\) by an unpebbled connector \(R\).  The
levels of the two endpoints differ, so \(|R|\ge2\).  Append \(Q\), \(R\), and the
old unpebbled path from \(y\) to \(S\).  Concatenation identifies two pairs of
endpoints, while \(R\) supplies at least two nodes; hence the result has length at least
\[
 \alpha_\pi n+z(\alpha_\pi-\sigma)n=p(z+1).
\]
For a node in the new source prefix, use its suffix of \(Q\) instead of all of
\(Q\).  That suffix contributes at least \((\alpha_\pi-\sigma)n\), and the old
path contributes \(z(\alpha_\pi-\sigma)n\).  This proves the stronger invariant
for the next chain source.  All pieces lie in footprints, hence are unpebbled;
acyclicity of the layered graph makes the concatenation a directed path rather than a
walk with repeated nodes.
\end{proof}

\begin{proof}[Proof of \cref{thm:latency}]
Let the constants be as in \eqref{eq:cap-search}--\eqref{eq:zmin}.  Footprint bounds
lower-bound the corresponding footprints, so every fertile challenge footprint has
weight at least \(\pi\).  Thus \eqref{eq:depth-robust} constructs the first chain source of
each segment.  The same applies to each source footprint, so every successful extension
constructs the next chain source.

When \(\ell>s_0\), \cref{lem:chain-length} gives a segment with
\(z\ge z_{\min}(\ell)\).
By \cref{lem:path-payoff}, that segment yields an unpebbled path ending in \(S\).
Since \(\sigma<\alpha_\pi\), the lower bound is increasing in \(z\), and hence
\[
 |Q|\ge
 \alpha_\pi n+
 \bigl(z_{\min}(\ell)-1\bigr)
 (\alpha_\pi-\sigma)n.
\]
Every piece used in the concatenation belongs to the footprint of \(S\).
Since \(s_0,s_1,h_0\), and \(h_1\) depend only on the scalar parameters,
\(z_{\min}(\ell)=\Omega(\ell)\), and the final claim follows.
\end{proof}

% References are included directly so this document has no BibTeX dependency.
\begin{thebibliography}{Rey24}

\bibitem[Fis19]{EC:Fisch19}
Ben Fisch.
\newblock Tight proofs of space and replication.
\newblock In Yuval Ishai and Vincent Rijmen, editors, {\em EUROCRYPT~2019,
  Part~II}, volume 11477 of {\em {LNCS}}, pages 324--348. Springer, Cham, May
  2019.

\bibitem[Rey24]{FOCS:Reyzin24}
Leonid Reyzin.
\newblock Proofs of space with maximal hardness.
\newblock In {\em 65th FOCS}, pages 1159--1177. {IEEE} Computer Society Press,
  October 2024.

\end{thebibliography}

\appendix
\crefalias{section}{appendix}
\crefalias{section}{appendix}
\crefalias{subsection}{appendix}
\renewcommand{\theHsection}{\Alph{section}}
\renewcommand{\theHsubsection}{\Alph{section}.\arabic{subsection}}

\section{Filecoin parameters}
\label{app:filecoin}
% Filecoin-shaped specialization.
Let \(\beta_8\) denote the degree-eight Chung expansion profile of
Reyzin~\cite[Claims 11 and 13]{FOCS:Reyzin24}.  We assume its concavity and
unique-maximizer condition.  The scalar parameters are the Filecoin-shaped SDR choice
of Fisch~\cite{EC:Fisch19}; a subscript \({\rm F}\) marks this specialization.

The chosen parameters are
\begin{equation}
 \delta_{\rm F}=0.0378,\quad
 \pi_{\rm F}=0.8,\quad
 \rho_{\rm F}=0.8,\quad
 \zeta_{\delta,{\rm F}}=0.8622,\quad
 \sigma_{\rm F}=0.1184,\quad
 \alpha_{\pi,{\rm F}}=0.2,\quad
 \sigmatilde_{\rm F}=3/5 .
 \label{eq:filecoin-parameters}
\end{equation}
The curve contributes \(\beta_8(\pi_{\rm F})=0.94911\ldots\) and
\(\beta_8(\sigma_{\rm F})=0.37888\ldots\), so the two zeros of \(\gaind\) are
\(\alpha_\delta^{\min}=0.00976\ldots\) and \(\alpha_\delta^{\max}=0.95243\ldots\), and
the expansion calculus of \cref{sec:model} takes the values of \cref{tab:filecoin}.
Every integer there is certified in the formal development from the rational bracket
\begin{equation}
 0.1113<g_{\pi,{\rm F}}<0.1114
 \label{eq:gpi-bracket}
\end{equation}
alone; no floating-point evaluation of \(\beta_8\) enters a proof.

\begin{table}[t]
\caption{The tracking constants and the capacities at the parameters
\eqref{eq:filecoin-parameters}.  The middle column names the definition being
evaluated.}
\label{tab:filecoin}
\begin{center}
\begin{tabular}{lll}
\toprule
Quantity & Definition & Value \\
\midrule
\(\bar\pi_{\rm F}\)          & \(1-\beta_8(\pi)\)                        & \(0.05088\ldots\) \\
\(g_{\pi,{\rm F}}\)          & \(\gaind(\pi)\)                           & \(0.11131\ldots\) \\
\(\gaind(\sigma_{\rm F})\)   & --                                        & \(0.22268\ldots\ge2g_{\pi,{\rm F}}\) \\
\(\pihat_{\rm F}\)           & \eqref{eq:tracking}, \(=\bar\pi_{\rm F}\) & \(0.05088\ldots\) \\
\(\ghat_{\rm F}\)            & \eqref{eq:tracking}, \(=g_{\pi,{\rm F}}\) & \(0.11131\ldots\) \\
\(\gtilde_{\rm F}\)          & \eqref{eq:gzeta}, \(=g_{\pi,{\rm F}}\)    & \(0.11131\ldots\) \\
\(\gaind(\sigmatilde_{\rm F})\) & --                                     & \(0.23406\ldots\ge2g_{\pi,{\rm F}}\) \\
\(\Pot_{\rm F}(\pi)\)        & \eqref{eq:midpoint}, \(=0.4408/g_{\pi,{\rm F}}\)         & \(3.9600\ldots\) \\
\midrule
\(s_{\rm F}\)                & \eqref{eq:cap-search}: \(\lceil6.628\ldots\rceil_++\lceil7.187\ldots\rceil-1\) & \(7+8-1=14\) \\
\(h_{0,{\rm F}}\)            & \eqref{eq:cap-span}: \(\lfloor13.310\ldots\rfloor+2\lceil7.187\ldots\rceil\) & \(13+16=29\) \\
\(h_{1,{\rm F}}\)            & \eqref{eq:cap-span-global}: \(1+\min\{6.1233\ldots,4.9600\ldots\}\) & \(5.9600\ldots\) \\
\(b^{\max}_{\rm F}\)         & \eqref{eq:cap-break}: \(\lceil0.929\ldots\rceil-1\)                          & \(0\) \\
\midrule
\(s_{0,{\rm F}}\)            & \(s_{\rm F}+b^{\max}_{\rm F}h_{0,{\rm F}}\)               & \(14\) \\
\(s_{1,{\rm F}}\)            & \(s_{\rm F}+1.6/g_{\pi,{\rm F}}\)                         & \(28.3740\ldots\) \\
\(s_{2,{\rm F}}\)            & \eqref{eq:joint-constants}: \((1-0.0622/g_{\pi,{\rm F}})+1.6/g_{\pi,{\rm F}}\) & \(0.4412\ldots+14.3740\ldots\) \\
\bottomrule
\end{tabular}
\end{center}
\end{table}

Beyond \eqref{eq:scalar-conditions}, the specialization verifies
\begin{equation}
 \bar\pi<\zeta_\delta-\rho,\qquad
 \zeta_\delta\le\alpha_\delta^{\max},\qquad
 \gaind(\sigma)\ge2\gpi,\qquad
 \rho<\beta_\delta(\pi)-\bar\pi,\qquad
 \gaind(\sigmatilde)\ge2\gpi ,
 \label{eq:no-break-conditions}
\end{equation}
the last of which is \eqref{eq:midpoint} at \(\ghat=\gpi\), with
\(\gaind(3/5)=0.23406\ldots\) against \(2g_{\pi,{\rm F}}=0.22262\ldots\).  Three of
these entries carry the whole specialization.  The condition
\(\gaind(\sigma)\ge2\gpi\) holds with little room,
\(\gaind(\sigma_{\rm F})=0.22268\ldots\) against
\(2g_{\pi,{\rm F}}=0.22262\ldots\), and it collapses \(\ghat_{\rm F}\) to
\(g_{\pi,{\rm F}}\).  The room is small because this condition is what fixes
\(\sigma_{\rm F}\) from below: the per-link payoff \((\alpha_\pi-\sigma)n\) rewards a
small \(\sigma\), while \(\gaind\) is still increasing there, so the least admissible
source weight is \(\sigma_{\min}=0.11834\ldots\) and the four-digit value \(0.118\) is
already inadmissible.  The challenge floor
\(\zeta_{\delta,{\rm F}}-\rho_{\rm F}=0.0622\) sits above \(\bar\pi_{\rm F}\) and has
gain \(0.13476\ldots>g_{\pi,{\rm F}}\), which collapses \(\gtilde_{\rm F}\) as well.
Finally, \(\beta_\delta(\pi_{\rm F})-\pihat_{\rm F}=0.86042\ldots\) exceeds the entire
budget \(\rho_{\rm F}=0.8\), so \(b^{\max}_{\rm F}=0\) and no break can be paid for.

\begin{corollary}[Filecoin-shaped specialization]
\label{cor:filecoin}
Assume inter-layer expansion according to \(\beta_8\) and
\eqref{eq:depth-robust} with \(\alpha_{\pi,{\rm F}}=0.2\).  Whenever \(\ell>14\), every
pebble configuration satisfying \eqref{eq:budgets} and every challenge set satisfying
\eqref{eq:challenge} have a footprint containing an unpebbled path of length at least
\[
 0.2\,n+\bigl(z_{\min,{\rm F}}(\ell)-1\bigr)0.0816\,n,
\]
where, writing \(h_{1,{\rm F}}=0.4408/g_{\pi,{\rm F}}+2\),
\begin{equation}
 z_{\min,{\rm F}}(\ell)=\max\left\{1,
 \left\lceil\frac{\ell-14-1.6/g_{\pi,{\rm F}}}{h_{1,{\rm F}}}\right\rceil_+,
 \left\lceil\frac{\ell-1+0.0622/g_{\pi,{\rm F}}-1.6/g_{\pi,{\rm F}}}
   {h_{1,{\rm F}}}\right\rceil_+,
 \left\lfloor\frac{\ell-14}{29}\right\rfloor+1\right\}.
 \label{eq:filecoin-zmin}
\end{equation}
\end{corollary}

\begin{proof}
The verified inequalities imply \eqref{eq:scalar-conditions}.  Apply
\cref{thm:latency}, whose hypothesis \(\sigma<\alpha_\pi\) reads
\(0.1184<0.2\).  Because \eqref{eq:no-break-conditions} is among what is verified
above, \(b^{\max}_{\rm F}=0\), so
\(s_{0,{\rm F}}=s_{\rm F}=14\) and \(s_{2,{\rm F}}\) is finite.
Substitution of \(h_{1,{\rm F}}\), \(s_{1,{\rm F}}\), \(s_{2,{\rm F}}\) and
\(h_{0,{\rm F}}=29\) in \eqref{eq:zmin} gives
\eqref{eq:filecoin-zmin}; also
\(\alpha_{\pi,{\rm F}}-\sigma_{\rm F}=0.0816\).
\end{proof}

The joint entry of \eqref{eq:filecoin-zmin} exceeds \(1\) exactly when
\(\ell>s_{2,{\rm F}}+h_{1,{\rm F}}=3+1.9786/g_{\pi,{\rm F}}\),
and the bracket \eqref{eq:gpi-bracket} places that quantity in \((20.761,20.778)\); the
plain ledger entry needs
\(\ell>s_{1,{\rm F}}+h_{1,{\rm F}}=16+2.0408/g_{\pi,{\rm F}}\in(34.319,34.337)\) and the
constant-charge entry \(\ell\ge43\).  All entries are nondecreasing in \(\ell\), so
\begin{equation}
 z_{\min,{\rm F}}(\ell)>1\iff \ell\ge21
 \qquad(\ell>14).
 \label{eq:filecoin-threshold}
\end{equation}
The smallest layer count for which the certified latency exceeds \(0.2n\) is therefore
\(\ell=21\), where the bound is \(0.2816n\).  The same bracket gives
\(z_{\min,{\rm F}}(1000)=166\), a bound of \(13.664n\).

\section{The Lean development}
\label{app:lean}
% Map from the explanatory note to Lean declarations.
The development is a single Lean~4 library; running \texttt{lake build} in the
repository's \texttt{lean/} directory rechecks it.  It uses no \texttt{sorry} and
declares no axioms of its own.
\Cref{tab:lean} names, for every numbered item of the paper, the Lean declaration that
states or proves it.

\begin{longtable}{@{}p{0.36\textwidth}p{0.60\textwidth}@{}}
\caption{Every numbered item of the paper and the Lean declaration that states or proves
it.  Module names are omitted; each declaration is unique in the library.}
\label{tab:lean}\\
\toprule
Paper & Lean \\
\midrule
\endfirsthead
\toprule
Paper & Lean \\
\midrule
\endhead
\multicolumn{2}{@{}l}{\emph{\Cref{sec:model}: model and main result}} \\
\eqref{eq:expansion}, expansion profile      & \texttt{LayeredGraph.expands} \\
\eqref{eq:depth-robust}, depth robustness    & \texttt{LayeredGraph.DepthRobust}; deletion
                                form \texttt{NodeDepthRobustAt}, equivalence
                                \texttt{depthRobustAt\_iff\_nodeDepthRobustAt} \\
\eqref{eq:budgets}, pebble budgets           & \texttt{Budget} \\
\eqref{eq:challenge}, challenge set          & \texttt{ChallengeBound} \\
\eqref{eq:tracking}, \(\pihat\) and \(\ghat\) & \texttt{Tracking.lam},
                                \texttt{Tracking.ghat} \\
\eqref{eq:scalar-conditions}, scalar conditions & \texttt{GeneralRegime} \\
\eqref{eq:gzeta}, \(\gtilde\)                & \texttt{Setting.gtilde} \\
\eqref{eq:cap-search}, search capacity \(s\) & \texttt{sCap}, \texttt{sCapOf},
                                \texttt{infertileCap}, \texttt{blockedCap} \\
\eqref{eq:cap-span}, attempt span \(h_0\)    & \texttt{h\textsubscript{0}},
                                \texttt{growthCap}, \texttt{spendCap} \\
\eqref{eq:cap-span-global}, global span \(h_1\) & \texttt{h\textsubscript{1}},
                                \texttt{asymptoticGrowth} \\
\eqref{eq:cap-break}, break capacity         & \texttt{bMax} \\
\eqref{eq:global-constants}, \(s_0\), \(s_1\) & \texttt{s\textsubscript{0}},
                                \texttt{s\textsubscript{1}}, \texttt{ledgerSlack} \\
\eqref{eq:zmin}, chain length                & \texttt{zMin}; at \(b^{\max}=0\),
                                \texttt{zMinNoBreak} \\
\eqref{eq:main-path-length}, path length     & \texttt{latencyLength},
                                \texttt{latencyLength\_general\_lower\_bound} \\
\Cref{thm:latency}            & \texttt{latency\_general} \\
\midrule
\multicolumn{2}{@{}l}{\emph{\Cref{sec:footprint-proof}: footprint growth and search}} \\
\eqref{eq:footprint-bound}, footprint bound  & \texttt{IsFootprintBound},
                                \texttt{footprintBound}; dominance
                                \texttt{Pebbling.footprintBound\_le} \\
\eqref{eq:expandable}, expandable level      & \texttt{Expandable} \\
\Cref{lem:challenge-floor}    & \texttt{zetaFloor\_le} (floor),
                                \texttt{gtilde\_le\_gainD\_of\_infertile} (rate),
                                \texttt{infertile\_card\_le\_gen} (capacity) \\
\Cref{lem:infertile-capacity} & \texttt{infertile\_budget} \\
\Cref{lem:first-source}       & \texttt{search\_bound}, \texttt{search\_reaches},
                                \texttt{basecase\_gen}; pebbling half
                                \texttt{challenge\_fertile} \\
\Cref{lem:mirror-floor}       & \texttt{mirror\_floor} \\
\Cref{lem:growth-window}      & \texttt{growth\_window} \\
\Cref{lem:fertile-continuation} & \texttt{fertile\_continuation\_gen} \\
\Cref{lem:break-charge}       & \texttt{break\_charge} \\
\Cref{lem:extension-attempt}  & \texttt{extension\_attempt\_gen} \\
\Cref{lem:infertile-window}   & \texttt{ChallengeBound.infertile\_card\_charge} \\
\eqref{eq:midpoint}, the mid-point \(\sigmatilde\) & \texttt{Tracking.mid}, \texttt{mid\_gain},
                                \texttt{two\_ghat\_le\_gainD\_of\_le\_mid} \\
\eqref{eq:midpoint}, \(\Pot\), and the growth term of \eqref{eq:cap-span-global} &
                                \texttt{growthPot}, \texttt{growthConst} \\
\Cref{lem:potential-window}   & \texttt{growthPot\_window};
                                \eqref{eq:growth-constant-window} is
                                \texttt{growthConst\_window} \\
\eqref{eq:potential-gap}      & \texttt{growthPot\_pi\_add\_gap};
                                the criterion is
                                \texttt{growthPot\_pi\_succ\_lt\_asymptoticGrowth} \\
\midrule
\multicolumn{2}{@{}l}{\emph{\Cref{sec:latency}: the latency theorem}} \\
\Cref{lem:global-ledger}      & \texttt{general\_ledger} (fused; see above) \\
\Cref{lem:joint-ledger}       & the \(nb=0\) clause of \texttt{GenLedger}, proved in
                                \texttt{general\_ledger}; read off by
                                \texttt{exists\_many\_links\_gen} \\
\eqref{eq:joint-constants}, \(s_2\) & \texttt{s\textsubscript{2}}, whose two
                                summands are \texttt{searchHead}
                                (\texttt{restartHead} at a general depth) and
                                \texttt{jointSlack}; \(\min\{\ghat,\gtilde\}\) is
                                \texttt{gmin} \\
third entry of \eqref{eq:zmin} & \texttt{jointEntry}, whose \texttt{if} on
                                \(b^{\max}=0\) is the \(s_2=\infty\) case;
                                \texttt{jointRatio\_le\_zMin} \\
\Cref{lem:chain-length}       & \texttt{exists\_many\_links\_gen} \\
\Cref{lem:path-payoff}        & \texttt{Link.base}, \texttt{Link.extend},
                                \texttt{chainSystem}, \texttt{chainPathLength} \\
\midrule
\multicolumn{2}{@{}l}{\emph{\Cref{app:filecoin}: Filecoin parameters}} \\
\eqref{eq:filecoin-parameters}, \cref{tab:filecoin}, the constants & \texttt{FilecoinLatencyParameters}; \texttt{sCap\_eq},
                                \texttt{h\textsubscript{0}\_eq}, \texttt{growthConst\_eq},
                                \texttt{h\textsubscript{1}\_eq},
                                \texttt{bMax\_eq\_zero}, \texttt{s\textsubscript{0}\_eq},
                                \texttt{gmin\_eq}, \texttt{searchHead\_eq},
                                \texttt{jointSlack\_eq} \\
\eqref{eq:gpi-bracket}, \(\gpi\) bracket     & \texttt{gpi8\_bounds} \\
\eqref{eq:no-break-conditions}, no-break conditions & \texttt{chung8GeneralRegime} \\
\Cref{cor:filecoin}           & \texttt{latency\_corollary}; on the constructed
                                degree-eight curve
                                \texttt{chung8\_latency\_corollary} \\
\eqref{eq:filecoin-zmin}, \(z_{\min,{\rm F}}\) & \texttt{filecoinZMin}, \texttt{zMin\_eq} \\
\eqref{eq:filecoin-threshold}, onset at \(\ell=21\) & \texttt{filecoinZMin\_20\_eq},
                                \texttt{filecoinZMin\_21\_eq},
                                \texttt{one\_lt\_filecoinZMin\_iff};
                                \texttt{latency\_21}; on the constructed curve
                                \texttt{chung8\_latency\_21} \\
\bottomrule
\end{longtable}

\end{document}
